\subsection{Diskretizimi i ekuacionit të valës}
Për
\begin{equation}
 u_{tt}=c^2u_{xx},
\end{equation}
vendosim $x_j=j\Delta x$ dhe $t^n=n\Delta t$. Diferencat qendrore japin
\begin{equation}
 \frac{u_j^{n+1}-2u_j^n+u_j^{n-1}}{\Delta t^2}
 =c^2\frac{u_{j+1}^n-2u_j^n+u_{j-1}^n}{\Delta x^2}
 +\Oh(\Delta t^2+\Delta x^2).
\end{equation}
Pas zgjidhjes për shtresën e re,
\begin{equation}
 u_j^{n+1}=2u_j^n-u_j^{n-1}
 +\lambda^2(u_{j+1}^n-2u_j^n+u_{j-1}^n),
 \qquad \lambda=\frac{c\Delta t}{\Delta x}.
\end{equation}
Skema është eksplicite dhe e rendit të dytë në hapësirë e kohë. Për të nisur, nga $u_t(x,0)=v_0(x)$ merret
\begin{equation}
 u_j^1=u_j^0+\Delta t\,v_{0,j}
 +\frac{\lambda^2}{2}(u_{j+1}^0-2u_j^0+u_{j-1}^0).
\end{equation}

\subsection{Analiza von Neumann dhe kushti CFL}
Provojmë një mod Fourier $u_j^n=G^ne^{\mathrm{i}j\theta}$. Zëvendësimi në skemë jep
\begin{equation}
 G^2-2\left(1-2\lambda^2\sin^2\frac\theta2\right)G+1=0.
\end{equation}
Qëndrueshmëria kërkon që të dy rrënjët të kenë modul jo më të madh se një. Kjo ndodh kur
\begin{equation}
 \lambda=\frac{c\Delta t}{\Delta x}\le1.
\end{equation}
Interpretimi është kauzal: gjatë një hapi kohor, informacioni fizik përshkon largësinë $c\Delta t$, e cila nuk duhet të tejkalojë shtrirjen $\Delta x$ që skema përdor për komunikim ndërmjet nyjeve.

\subsection{Dispersimi numerik}
Duke vendosur $G=e^{-\mathrm{i}\omega\Delta t}$ merret lidhja e dispersimit
\begin{equation}
 \sin^2\left(\frac{\omega\Delta t}{2}\right)
 =\lambda^2\sin^2\left(\frac{k\Delta x}{2}\right).
\end{equation}
Në vazhdësi $\omega=ck$, por në rrjet raporti $\omega/k$ varet nga $k$. Modet me gjatësi vale të krahasueshme me $\Delta x$ përhapen me shpejtësi fazore të gabuar. Prandaj kushti CFL garanton qëndrueshmëri, jo saktësi. Një rregull praktik është të përdoren shumë pika për gjatësi vale dhe të kryhet studim rrjeti.

\subsection{Kushtet kufitare dhe energjia}
Kufiri i fiksuar përdor $u=0$ dhe reflekton me ndryshim shenje. Kufiri i lirë përdor $u_x=0$. Një kusht dalës i rendit të parë,
\begin{equation}
 u_t+cu_x=0,
\end{equation}
lejon kalimin e valës djathtas, por diskretizimi i tij mund të reflektojë pjesërisht. Energjia e vazhdueshme
\begin{equation}
 E(t)=\frac12\int\left(u_t^2+c^2u_x^2\right)\dd x
\end{equation}
është konstante për kufij pa fluks. Një energji diskrete e qëndrueshme është provë e fortë e implementimit.

\subsection{Valët dhe reflektimi nga plazma}
Zgjidhja e d'Alembert-it $u(x,t)=F(x-ct)+G(x+ct)$ ndan valët bredhëse djathtas e majtas. Në një interval me skaje të fiksuara, kushtet kufitare zgjedhin numra valorë diskretë $k_n=n\pi/L$ dhe frekuenca $\omega_n=ck_n$, duke formuar valë të qëndrueshme.

Në një plazmë të ftohtë pa fushë magnetike, permitiviteti relativ është
\begin{equation}
 \varepsilon_r(\omega)=1-\frac{\omega_p^2}{\omega^2},\qquad
 \omega_p=\sqrt{\frac{n_ee^2}{m_e\varepsilon_0}}.
\end{equation}
Për $\omega<\omega_p$, numri valor bëhet imagjinar dhe fusha shuhet brenda plazmës; vala reflektohet. Në modelin me përplasje, $\omega^2$ zëvendësohet në mënyrë efektive nga $\omega(\omega+\mathrm{i}\nu)$ dhe shfaqet absorbimi. Simulimi duhet të dallojë reflektimin fizik nga ai artificial i kufirit numerik.

\subsection{Përmbledhje dhe ushtrime}
Diskretizimi i valëve kërkon njëkohësisht konsistencë, CFL, rezolucion spektral dhe kushte kufitare fizikisht të përshtatshme.
\begin{enumerate}
\item Nxirrni formulën e shtresës së parë nga Taylor-i dhe shpejtësia fillestare.
\item Provoni kushtin CFL me analizën von Neumann.
\item Matni shpejtësinë fazore numerike për disa numra valorë.
\item Ndërtoni një valë të qëndrueshme dhe verifikoni frekuencat normale.
\item Simuloni kalimin nga vakumi në plazmë dhe ndani reflektimin fizik nga ai i kufirit të djathtë.
\end{enumerate}
